% -------------------------------------- Modèles : Partie photovoltaïque ------------------------------------------
\begin{frame}[t]{Modèles : partie photovoltaïque}

  \vspace{-1em}

  \begin{columns}[T,onlytextwidth]

    % ----------- Schéma -----------
    \column{0.48\textwidth}
\begin{pccard}
  \centering
  \figschemaphotovoltaique
  {\scriptsize\bfseries\color{pcNavy}
   Représentation de la partie photovoltaïque}
\end{pccard}


\begin{pcbanner}{Démarche}
  \begin{itemize}\footnotesize
    \item Résolution 1D du système couplé
    \item Comparer aux résultats de SCAPS-1D
    \item Confronter à l'expérimental (ANR)
  \end{itemize}
\end{pcbanner}

    % ----------- Approches -----------
    
        \column{0.48\textwidth}
\vspace{-0.5em}
\begin{pcbanner}{Approche cinétique}
  \vspace{0.45em}
  \begin{pcbannerlight}{Principe}
    \vspace{0.3em}
    \begin{itemize}\footnotesize
      \item Suivi local des porteurs et du potentiel électrostatique
      \item Poisson + équations de continuité
      \item Transport par dérive--diffusion
    \end{itemize}
  \end{pcbannerlight}

  \vspace{0.35em}

  \begin{pcbannerlight}{Physique prise en compte}
    \vspace{0.3em}
    \begin{itemize}\footnotesize
      \item Photogénération (électromagnétisme ou radiatif)
      \item Recombinaisons : SRH, Auger, radiative
      \item Hétérojonction CIGS/CdS/ZnO
    \end{itemize}
  \end{pcbannerlight}
  \vspace{0.2em}
\end{pcbanner}

  \end{columns}


\end{frame}

%------------------------------------------------------------------------

\begin{frame}[t]{Modèles : partie photovoltaïque}

  \vspace{-1em}

  \begin{columns}[T,onlytextwidth]

    % ----------- Schéma -----------
    \column{0.48\textwidth}
\begin{pccard}
  \centering
  \figschemaphotovoltaique

  {\scriptsize\bfseries\color{pcNavy}
   Représentation de la partie photovoltaïque}
\end{pccard}


\begin{pcbanner}{Photogénération}
      \centering\scriptsize
      \textbf{Photogénération (Bouguer--Lambert--Beer)}
      \[ G(x) = q_0\,k_a\,e^{-k_a x} \]
      \end{pcbanner}

    % ----------- Approches -----------
    
    \column{0.48\textwidth}
\vspace{-0.5em}
    \begin{pcbanner}{Système de Van Roosbroeck}
      \vspace{0.3em}
      \centering\scriptsize

      \textbf{Équation de Poisson}
      \[ -\nabla\!\cdot\!\left(\varepsilon\,\nabla\psi\right) = q\,(C + p - n) \]\\

\vspace{1em}
      \textbf{Équations de continuité}
      \[ \frac{\partial n}{\partial t} = \frac{1}{q}\,\nabla\!\cdot\!\vec{J}_n - R + G_{th} + G_{ph} \]
      \[ \frac{\partial p}{\partial t} = -\frac{1}{q}\,\nabla\!\cdot\!\vec{J}_p - R + G_{th} + G_{ph} \]\\

\vspace{1em}
      \textbf{Densités de courant (dérive-diffusion)}
      \[ \vec{J}_n = q\,\mu_n\!\left(-n\,\nabla\psi + U_t\,\nabla n\right) \]
      \[ \vec{J}_p = q\,\mu_p\!\left(-p\,\nabla\psi - U_t\,\nabla p\right) \]\\
\vspace{1em}
      \textbf{CL Robin ou Ohmique}
       \vspace{-0.5em}
    \end{pcbanner}

  \end{columns}


\end{frame}